% !TEX root = print.tex
% This article's additional theorem-like environments.
%
% Framework-owned theorems.tex defines definition/lemma/proposition/theorem on one
% counter mastered by `definition` and reset per chapter. Three more are needed here,
% and all three ride the SAME counter, because the numbering scheme below depends on it.
%
% Why they are not in macros.tex: both drivers (print.tex, web.tex) \input macros
% *before* theorems, so `definition` does not exist yet at that point. This file is
% \input from the top of content.tex instead, after theorems.tex has run.
%
%   corollary -- omitted by theorems.tex although `linkage` counts it a statement
%                environment (config.py DEFAULT_STATEMENT_ENVS) and plastexdepgraph
%                graphs it by default. The draft has three. A `corollary` here IS a
%                blueprint node and needs \label/\statusT/\statusA like any other.
%
%   remark    -- NOT a statement environment: no \label discipline, no status tag, not
%   example      in the dependency graph. They carry the draft's Remarks and Examples so
%                the blueprint stays the full text of record. A remark or example whose
%                content is load-bearing must be promoted to a definition/proposition
%                instead -- see the \S8 families, where the draft's Examples 8.1-8.3
%                become nodes.
%
% Adding a *statement* kind beyond corollary would also mean listing it in web.tex's
% `thms=` option, or the graph would silently ignore its \uses/\leanok annotations.

\newtheorem{corollary}[definition]{Corollary}
\newtheorem{remark}[definition]{Remark}
\newtheorem{example}[definition]{Example}
